\documentclass[border=12pt]{standalone}
\usepackage[no-math]{fontspec}
\usepackage[brazil]{babel}
\usepackage{amsmath,amssymb}
\IfFileExists{C:/Windows/Fonts/Lato_Regular.ttf}{%
  \setmainfont{Lato}[
    Path = C:/Windows/Fonts/,
    UprightFont = Lato_Regular.ttf,
    BoldFont = Lato_Bold.ttf,
    ItalicFont = Lato_Italic.ttf,
    BoldItalicFont = Lato_BoldItalic.ttf
  ]
  \setsansfont{Lato}[
    Path = C:/Windows/Fonts/,
    UprightFont = Lato_Regular.ttf,
    BoldFont = Lato_Bold.ttf,
    ItalicFont = Lato_Italic.ttf,
    BoldItalicFont = Lato_BoldItalic.ttf
  ]
}{%
  \RequirePackage[default]{lato}
}
\renewcommand{\familydefault}{\sfdefault}

\usepackage{xcolor}
\usepackage{tikz}
\usepackage{pgfplots}
\pgfplotsset{compat=1.18}
\usetikzlibrary{arrows.meta,positioning,shapes.geometric,fit,backgrounds,calc,decorations.pathreplacing,patterns}

% Paleta institucional UEFS / Extensão Rural
\definecolor{erNavy}{HTML}{1F4268}
\definecolor{erNavyDark}{HTML}{152E4A}
\definecolor{erGold}{HTML}{B38A3D}
\definecolor{erGoldDark}{HTML}{8C6927}
\definecolor{erTeal}{HTML}{0E7C7B}
\definecolor{erForest}{HTML}{2D6A4F}
\definecolor{erAlert}{HTML}{9E2A2B}
\definecolor{erInk}{HTML}{26313D}
\definecolor{erMuted}{HTML}{5B6573}

\definecolor{bgNavy}{HTML}{EDF2F7}
\definecolor{bgGold}{HTML}{F8F3E9}
\definecolor{bgTeal}{HTML}{EBF6F5}
\definecolor{bgForest}{HTML}{EEF5F0}
\definecolor{bgAlert}{HTML}{FDF2F2}

\tikzset{
  >=Stealth,
  erBox/.style={
    draw=#1, fill=white, rounded corners=5pt, line width=1.15pt,
    inner sep=7pt, align=center, font=\footnotesize\sffamily, text=erInk
  },
  erTint/.style={
    draw=#1, fill=#1!8!white, rounded corners=5pt, line width=1.15pt,
    inner sep=7pt, align=center, font=\footnotesize\sffamily, text=erInk
  },
  erBadge/.style={
    fill=#1, text=white, rounded corners=9pt, inner xsep=8pt, inner ysep=3pt,
    font=\scriptsize\bfseries\sffamily
  },
  erTitle/.style={font=\normalsize\bfseries\sffamily, text=erNavy, align=center},
  erSub/.style={font=\footnotesize\sffamily, text=erMuted, align=center},
  erArrow/.style={->, >=Stealth, line width=1.35pt},
  erDArrow/.style={<->, >=Stealth, line width=1.35pt}
}


\begin{document}
\begin{tikzpicture}

  % Título geral no topo
  \node[erTitle, font=\large\bfseries\sffamily] (title) at (7.65, 8.05)
    {Difusão de Inovações (Rogers) e Estratégia de ATER Digital Híbrida};

  % ================= ESQUERDA: PGFPLOTS CURVA S =================
  \begin{scope}[shift={(0, 0)}]
    % Moldura sutil para o painel esquerdo (alinhada ao painel direito)
    \draw[draw=erNavy!30, fill=bgNavy!45, rounded corners=8pt, line width=1pt]
      (-1.35, -1.75) rectangle (8.35, 7.35);

    \node[erBadge=erNavy, rounded corners=5pt, inner ysep=4pt, anchor=west] at (-1.05, 6.85)
      {Curva S de Difusão de Inovações (Rogers, 2003)};

    \begin{axis}[
      at={(-0.15cm, -0.35cm)},
      width=8.1cm,
      height=5.95cm,
      scale only axis,
      xmin=0, xmax=10,
      ymin=0, ymax=106,
      xtick={1.67, 3.49, 5.0, 6.51},
      xticklabels={},
      ytick={0, 25, 50, 75, 100},
      yticklabels={0\%, 25\%, 50\%, 75\%, 100\%},
      xlabel={Tempo de Adoção},
      ylabel={\% de Adoção Acumulada},
      label style={font=\footnotesize\bfseries\sffamily, text=erInk},
      tick label style={font=\scriptsize\sffamily, text=erMuted},
      axis lines=left,
      axis line style={draw=erNavy, line width=1pt},
      ymajorgrids=true,
      grid style={dashed, draw=erNavy!18},
      clip=false
    ]
      % Linhas verticais tracejadas separando os 5 estágios acumulados (2.5%, 16%, 50%, 84%)
      \draw[dashed, color=erTeal!55, line width=0.8pt] (axis cs:1.67, 0) -- (axis cs:1.67, 100);
      \draw[dashed, color=erTeal!55, line width=0.8pt] (axis cs:3.49, 0) -- (axis cs:3.49, 100);
      \draw[dashed, color=erTeal!55, line width=0.8pt] (axis cs:5.00, 0) -- (axis cs:5.00, 100);
      \draw[dashed, color=erTeal!55, line width=0.8pt] (axis cs:6.51, 0) -- (axis cs:6.51, 100);

      % Curva logística S suave: 100 / (1 + exp(-1.1*(x - 5)))
      \addplot[
        domain=0.2:9.8,
        samples=120,
        smooth,
        line width=2.4pt,
        color=erNavy
      ] {100 / (1 + exp(-1.1*(x - 5)))};

      % Pontos marcadores nos estágios
      \addplot[
        only marks,
        mark=*,
        mark size=2.8pt,
        color=erGoldDark,
        fill=erGold
      ] coordinates {
        (0.95, 1.15)
        (2.58, 6.53)
        (4.25, 30.5)
        (5.75, 69.5)
        (7.75, 95.4)
      };

      % Rótulos sem sobreposição com linhas conectoras limpas
      \node[erBox=erGold, inner sep=3.5pt, font=\scriptsize\sffamily, fill=white]
        (lbl1) at (axis cs:1.25, 22) {\textbf{Inovadores}\\2,5\%};
      \draw[->, >=Stealth, color=erGoldDark, line width=0.8pt] (lbl1.south) -- (axis cs:0.95, 3.5);

      \node[erBox=erTeal, inner sep=3.5pt, font=\scriptsize\sffamily, fill=white]
        (lbl2) at (axis cs:2.35, 45) {\textbf{Adotantes}\\[-1pt]\textbf{Iniciais}\\13,5\%};
      \draw[->, >=Stealth, color=erTeal, line width=0.8pt] (lbl2.south) -- (axis cs:2.58, 9.2);

      \node[erBox=erForest, inner sep=3.5pt, font=\scriptsize\sffamily, fill=white]
        (lbl3) at (axis cs:3.25, 76) {\textbf{Maioria}\\[-1pt]\textbf{Inicial}\\34\%};
      \draw[->, >=Stealth, color=erForest, line width=0.8pt] (lbl3.south east) -- (axis cs:4.15, 33.5);

      \node[erBox=erNavy, inner sep=3.5pt, font=\scriptsize\sffamily, fill=white]
        (lbl4) at (axis cs:7.6, 52) {\textbf{Maioria}\\[-1pt]\textbf{Tardia}\\34\%};
      \draw[->, >=Stealth, color=erNavy, line width=0.8pt] (lbl4.north west) -- (axis cs:5.9, 67.5);

      \node[erBox=erMuted, inner sep=3.5pt, font=\scriptsize\sffamily, fill=white]
        (lbl5) at (axis cs:8.35, 81) {\textbf{Retardatários}\\16\%};
      \draw[->, >=Stealth, color=erMuted, line width=0.8pt] (lbl5.north) -- (axis cs:7.75, 93);

    \end{axis}
  \end{scope}

  % ================= DIREITA: ATER DIGITAL HÍBRIDA =================
  \begin{scope}[shift={(12.5, 0)}]
    % Moldura externa do painel direito (mesma altura exata do painel esquerdo)
    \draw[draw=erTeal!60, fill=bgTeal!70, rounded corners=8pt, line width=1.15pt]
      (-3.7, -1.75) rectangle (3.7, 7.35);

    \node[erBadge=erTeal, rounded corners=5pt, inner ysep=4pt] (rbadge) at (0, 6.85)
      {ATER Digital Híbrida};
    \node[font=\scriptsize\bfseries\sffamily, text=erNavy, below=4pt of rbadge] (rsub)
      {Evitando a Exclusão Rural};

    % Bloco superior: Canais Digitais (posicionado com below=of para zero sobreposição)
    \node[erBox=erTeal, text width=6.3cm, align=left, inner sep=8pt, below=8pt of rsub]
      (dig) {
      {\color{erTeal}\textbf{\footnotesize Canais Digitais (Escala \& Agilidade)}}\\[4pt]
      \scriptsize
      \textbullet\ \textbf{WhatsApp \& Áudio:} alertas técnicos e dúvidas\\
      \textbullet\ \textbf{Rádio Web \& Podcasts:} alcance sem banda larga\\
      \textbullet\ \textbf{Sensores \& Plataformas:} clima, manejo e mercados
    };

    % Conector central
    \node[fill=erGold, text=white, circle, inner sep=3.5pt, font=\small\bfseries\sffamily,
          below=11pt of dig] (plus) {+};
    \draw[erArrow, color=erTeal] (dig.south) -- (plus.north);

    % Bloco inferior: Mediação Presencial
    \node[erBox=erNavy, text width=6.3cm, align=left, inner sep=8pt, below=11pt of plus]
      (pres) {
      {\color{erNavy}\textbf{\footnotesize Mediação Presencial (Vínculo \& Inclusão)}}\\[4pt]
      \scriptsize
      \textbullet\ \textbf{Visitas Técnicas:} diagnóstico situado na propriedade\\
      \textbullet\ \textbf{Unidades Demonstrativas:} aprendizado entre pares\\
      \textbullet\ \textbf{Acesso \& Conectividade:} inclusão e letramento digital
    };
    \draw[erArrow, color=erNavy] (plus.south) -- (pres.north);

    % Nota de rodapé do painel direito
    \node[fill=bgNavy, draw=erNavy!45, rounded corners=5pt, inner sep=6pt,
          text width=6.3cm, align=center, font=\scriptsize\bfseries\itshape\sffamily, text=erNavy,
          below=10pt of pres]
      {O digital potencializa a ATER sem substituir o contato humano.};
  \end{scope}

\end{tikzpicture}
\end{document}
